% Q17 -- Parity generator 1-bit cell, CMOS (schematic).
% Pull-down n-block shown explicitly (two series branches); complementary pull-up
% p-block shown as its dual network box; output buffered by a static CMOS inverter.
\documentclass[border=10pt]{standalone}
\usepackage{circuitikz}
\begin{document}
\begin{circuitikz}
\draw (-2,9) -- (6,9) node[right]{$V_{DD}$};
\draw (-2,0) -- (6,0);
\path (2,0) node[ground]{};
% pull-up p-block (dual network) as a box
\draw[thick] (-0.9,6.6) rectangle (0.9,8.2);
\node at (0,7.4) {p-block};
\draw (0,9) -- (0,8.2);
\draw (0,6.6) -- (0,5.6) coordinate (PB);
\path (PB) node[circ]{};
\node[above right] at (PB) {$\overline{P}_i$};
% n-block bus
\draw (PB) -- (0,5.0);
\draw (-0.8,5.0) -- (0.8,5.0);
% branch 1
\path (-0.8,3.6) node[nmos] (T1a){}; \draw (T1a.D) -- (-0.8,5.0);
\draw (T1a.G) -- (-1.9,3.6) node[left]{$A_i$};
\path (-0.8,1.8) node[nmos] (T1b){}; \draw (T1b.D) -- (T1a.S);
\draw (T1b.G) -- (-1.9,1.8) node[left]{$\overline{P}_{i-1}$};
\draw (T1b.S) -- (-0.8,0);
% branch 2
\path (0.8,3.6) node[nmos, xscale=-1] (T2a){}; \draw (T2a.D) -- (0.8,5.0);
\draw (T2a.G) -- (1.9,3.6) node[right]{$\overline{A}_i$};
\path (0.8,1.8) node[nmos, xscale=-1] (T2b){}; \draw (T2b.D) -- (T2a.S);
\draw (T2b.G) -- (1.9,1.8) node[right]{$P_{i-1}$};
\draw (T2b.S) -- (0.8,0);
% static CMOS inverter -> P_i
\path (4,7) node[pmos] (IP){}; \draw (IP.S) -- (4,9);
\path (4,2) node[nmos] (IN){}; \draw (IN.S) -- (4,0);
\draw (IP.D) -- (4,4.5) coordinate (PO);
\draw (IN.D) -- (PO);
\path (PO) node[circ]{};
\draw (PO) -- (6,4.5) node[right]{$P_i$};
\draw (IP.G) -- (2.6,7); \draw (IN.G) -- (2.6,2);
\draw (2.6,7) -- (2.6,2); \draw (2.6,4.5) -- (PB);
\end{circuitikz}
\end{document}
